\subsection{DERIVATIVES OF THE EXPONENTIAL FUNCTIONS; THE NUMBER e}

We know that every continuous homomorphism of the additive group R into the multiplicative group \(R^{*}\) of real numbers \(\neq 0\) is a function of the form \(x \mapsto a^{x}\) (called an exponential function) where a is a number \textgreater{} 0 (TG, V, p.11); it is an isomorphism of R onto the multiplicative group \(R_{+}^{*}\) of numbers \textgreater{} 0 if \(a \neq 1\) , and the inverse isomorphism from \(R_{+}^{*}\) onto R is denoted by \(\log_{a} x\) and is called the logarithm to the base a.

We shall see that the function \(f(x) = a^{x}\) has, for every \(x \in R\) , a derivative of the form \(c.a^{x}\) (where clearly \(c = f'(0)\) ). This results from the following general theorem:

THEOREM 1. Let E be a complete normed algebra over the field R, with a unit element e, and let f be a continuous group homomorphism of the additive group R into the multiplicative group G of invertible elements of E. Then the map f is differentiable at every \(x \in R\) , and

\[
\mathbf {f} ^ {\prime} (x) = \mathbf {f} (x) \mathbf {f} ^ {\prime} (0).\tag{1}
\]

First we note that, E being a complete algebra, G is open in E (Gen.~Top., IX, p.~179, prop. 14). Consider the function \(\mathbf{g}(x) = \int_0^a \mathbf{f}(x + t) dt\) , where \(a > 0\) is a number which we shall choose later; since \(\mathbf{f}(x + t) = \mathbf{f}(x) \mathbf{f}(t)\) by hypothesis, we have \(\mathbf{g}(x) = \int_0^a \mathbf{f}(x) \mathbf{f}(t) dt = \mathbf{f}(x) \int_0^a \mathbf{f}(t) dt\) (I, p.~6, prop. 3). Let \(\alpha > 0\) be such that the ball \(\| \mathbf{x} - \mathbf{e} \| \leqslant \alpha\) is contained in G; since \(\mathbf{f}(0) = \mathbf{e}\) and \(\mathbf{f}\) is continuous by hypothesis, one can assume that \(a\) is small enough so that \(\| \mathbf{f}(t) - \mathbf{e} \| \leqslant \alpha\) on {[}0, a{]}; consequently (II, p.~61, formula (16)) one has

\[
\left\| \frac {1}{a} \int_ {0} ^ {a} \mathbf {f} (t) d t - \mathbf {e} \right\| \leqslant \alpha ,
\]

and \(\frac{1}{a}\int_{0}^{a}\mathbf{f}(t)dt\) belongs to G; in other words, is invertible; so too is \(\mathbf{b}=\int_{0}^{a}\mathbf{f}(t)dt\) and one can write \(\mathbf{f}(x)=\mathbf{g}(x)\mathbf{b}^{-1}\) ; it is therefore enough to show that g is differentiable; now, by the change of variable \(x+t=u\) we have \(\mathbf{g}(x)=\int_{x}^{x+a}\mathbf{f}(u)du\) ; since f is continuous, g is differentiable for all \(x\in R\) (II, p.~56, prop. 3), and

\[
\mathbf {g} ^ {\prime} (x) = \mathbf {f} (x + a) - \mathbf {f} (x) = \mathbf {f} (x) (\mathbf {f} (a) - \mathbf {e}).
\]

Hence \(\mathbf{f}'(x) = \mathbf{g}'(x)\mathbf{b}^{-1} = \mathbf{f}(x)\mathbf{c}\) , where \(\mathbf{c} = (\mathbf{f}(a) - \mathbf{e})\mathbf{b}^{-1}\) , and clearly \(\mathbf{f}'(0) = \mathbf{c}\) .

Conversely, one can show, either directly (III, p.~115, exerc. 1), or by means of the theory of linear differential equations (IV, p.~188), that every differentiable map \(\mathbf{f}\) of \(\mathbf{R}\) into a complete normed algebra E, such that \(\mathbf{f}'(x) = \mathbf{f}(x)\mathbf{c}\) and \(\mathbf{f}(0) = \mathbf{e}\) , is a homomorphism of the additive group \(\mathbf{R}\) into the multiplicative group G.

PROPOSITION 1. For every number a \textgreater{} 0 and \(\neq 1\) the exponential function \(a^{x}\) admits at every point \(x \in R\) a derivative equal to \((\log_{e}a)a^{x}\) where e is a number \textgreater{} 1 (independent of a).

Applying th. 1 to the case where E is the field R itself now shows that \(a^{x}\) has a derivative equal to \(\varphi(a).a^{x}\) at every point, where \(\varphi(a)\) is a real number \(\neq 0\) depending only on a. Let b be a second number \textgreater{} 0 and \(\neq 1\) ; the function \(b^{x}\) has a derivative equal to \(\varphi(b).b^{x}\) from the above; on the other hand, we have \(b^{x} = a^{x.\log_{a}b}\) so (I, p.~9, prop. 5) the derivative of \(b^{x}\) is equal to \(\log_{a}b.\varphi(a)b^{x}\) ; on comparing these two expressions we obtain

\[
\varphi (b) = \varphi (a). \log_ {a} b.\tag{2}
\]

One deduces that there is one unique number \(b\) such that \(\varphi(b) = 1\) ; by (2) this relation is equivalent to \(b = a^{1/\varphi(a)}\) . It is conventional to denote the real number so obtained by \(e\) ; from (2) one has \(\varphi(a) = \log_e a\) , which completes the proof of prop. 1.

One often writes \(\exp x\) instead of \(e^x\) .

The definition of the number e shows that

\[
\mathrm{D} (e ^ {x}) = e ^ {x}\tag{3}
\]

which proves that \(e^{x}\) is strictly increasing, hence that e \textgreater{} 1.

In §2 (III, p.~105) we shall see how to calculate arbitrarily close approximations to \(e\) .

DEFINITION 1. Logarithms to the base \(e\) are called Naperian logarithms (or natural logarithms).

We usually omit the base in the notation for the Naperian logarithm. Unless it is stated to the contrary, the notation \(\log x (x > 0)\) will denote the Naperian logarithm of x. With this notation, prop. 1 can be written as the identity

\[
\mathrm{D} \left(a ^ {x}\right) = (\log a) a ^ {x}\tag{4}
\]

valid for any a \textgreater{} 0 ( \(\log a = 0\) when a = 1).

This relation shows that \(a^x\) has derivatives of every order, and that

\[
\mathrm{D} ^ {n} \left(a ^ {x}\right) = (\log a) ^ {n} a ^ {x}.\tag{5}
\]

In particular, for a \textgreater{} 0 and \(\neq 1\) one has \(\mathrm{D}^{2}(a^{x}) > 0\) for all \(x \in R\) , and hence \(a^{x}\) is strictly convex on R (I, p.~31, corollary). From this one deduces the following proposition:

PROPOSITION 2 (``geometric mean inequality''). For any numbers \(z_{i} > 0\) ( \(1 \leqslant i \leqslant n\) ) and numbers \(p_{i} > 0\) such that \(\sum_{i=1}^{n} p_{i} = 1\) , one has

\[
z _ {1} ^ {p _ {1}} z _ {2} ^ {p _ {2}} \dots z _ {n} ^ {p _ {n}} \leqslant p _ {1} z _ {1} + p _ {2} z _ {2} + \dots + p _ {n} z _ {n}.\tag{6}
\]

Moreover, the two sides of (6) are equal only if the \(z_{i}\) are equal.

Let us put \(z_{i} = e^{x_{i}}\) ; then the inequality (6) can be written

\[
\exp \left(p _ {1} x _ {1} + p _ {2} x _ {2} + \dots + p _ {n} x _ {n}\right) \leqslant p _ {1} e ^ {x _ {1}} + p _ {2} e ^ {x _ {2}} + \dots + p _ {n} e ^ {x _ {n}}.\tag{7}
\]

The proposition thus follows from prop. 1 of I, p.~26 applied to the function \(e^x\) , which is strictly convex on \(\mathbf{R}\) .

One says that the left- (resp. right-) hand side of (6) is the weighted geometric mean (resp. weighted arithmetic mean) of the n numbers \(z_{i}\) relative to the weights \(p_{i}\) ( \(1 \leqslant i \leqslant n\) ). If \(p_{i} = 1/n\) for \(1 \leqslant i \leqslant n\) , one calls the corresponding arithmetic and geometric means the ordinary arithmetic and geometric means of the \(z_{i}\) . Then the inequality (6) can be written

\[
\left(z _ {1} z _ {2} \dots z _ {n}\right) ^ {1 / n} \leqslant \frac {1}{n} \left(z _ {1} + z _ {2} + \dots + z _ {n}\right).\tag{8}
\]

